\documentclass[12pt,aspectratio=169]{beamer}
\usepackage[T1]{fontenc}
\usepackage{lmodern}
\usepackage{amsmath,amssymb,array,booktabs,tabularx,tikz}
% Add language support appropriate to the source, if needed.
% Keep this preamble inline in every delivered deck.
\definecolor{GEInk}{HTML}{132A3E}
\definecolor{GETeal}{HTML}{17675E}
\definecolor{GEFooter}{HTML}{EAF2F1}
\setbeamercolor{background canvas}{bg=white}
\setbeamercolor{normal text}{fg=GEInk,bg=white}
\setbeamercolor{structure}{fg=GETeal}
\setbeamercolor{frametitle}{fg=GEInk}
\setbeamercolor{GE footer}{fg=GETeal,bg=GEFooter}
\setbeamerfont{normal text}{family=\sffamily,size=\normalsize}
\setbeamerfont{title}{size=\huge,series=\bfseries}
\setbeamerfont{subtitle}{size=\normalsize,series=\bfseries}
\setbeamerfont{frametitle}{size=\Large,series=\bfseries}
\setbeamerfont{framesubtitle}{size=\normalsize,series=\mdseries}
\setbeamerfont{GE footer}{size=\footnotesize}
\setbeamersize{text margin left=9mm,text margin right=9mm}
\setbeamertemplate{navigation symbols}{}
\setbeamertemplate{headline}{}
\setbeamercovered{invisible}
\setbeamertemplate{frametitle}{%
  \vskip.5em
  {\usebeamerfont{frametitle}\strut\insertframetitle\strut\par}%
  \ifx\insertframesubtitle\empty\else
    {\usebeamerfont{framesubtitle}\strut\insertframesubtitle\strut\par}%
  \fi
  \vskip.25em
}
\setbeamertemplate{footline}{%
  \begin{beamercolorbox}[wd=\paperwidth,ht=3ex,dp=1.3ex,
    leftskip=9mm,rightskip=9mm]{GE footer}%
    \usebeamerfont{GE footer}\insertsectionhead\hfill
    \insertframenumber/\inserttotalframenumber
  \end{beamercolorbox}%
}
\setlength{\parindent}{0pt}
\setlength{\parskip}{.35em}
\setlength{\tabcolsep}{4pt}
\renewcommand{\arraystretch}{1.22}
\AtBeginEnvironment{frame}{\raggedright}
\addtobeamertemplate{frametitle}{}{%
  \setlength{\abovedisplayskip}{6pt plus 2pt minus 2pt}%
  \setlength{\belowdisplayskip}{6pt plus 2pt minus 2pt}%
  \setlength{\abovedisplayshortskip}{6pt plus 2pt minus 2pt}%
  \setlength{\belowdisplayshortskip}{6pt plus 2pt minus 2pt}%
}


% SOURCE MAP (source practice is read only):
% book2026-2027/grade11/term_2/1-practice_activities/G11_T2_L1_C2_practice.tex
% Introduction: sec:introduction through immediately before
% sec:notation-and-procedure; includes both subsections and its legend.
% Metadata: title and Evaluated criteria, C2 (verbatim below).
% Selected example: prob:c2-2 only; solutions sol:c2-2-a/b/c.
% Inventory: general prob:c2-1, c2-2, c2-3; advanced c2-4, c2-5;
% all have three linked solutions. C2-2 chosen for full comparison/rounding.
% Frames 1 metadata; 2/25 intro guide; 3 cycling motivation;
% 4 key terms; 5/6 guide's illustrative quantities; 7/8 cycling tasks 1/2;
% 9/10 cycling task 3 and interpretation; 11-23 prob:c2-2 + a/b/c;
% 24 summary from intro + c2-2 only. No notation/procedure content imported.
% The source-local orange palette overrides the registry's default palette.
\definecolor{PopulationMean}{HTML}{FF765C}
\definecolor{SampleMean}{HTML}{9BDD36}
\definecolor{PopulationSD}{HTML}{E42134}
\definecolor{SampleSize}{HTML}{173B17}
\definecolor{Confidence}{HTML}{55C8F0}
\definecolor{Critical}{HTML}{2648BB}
\definecolor{StandardError}{HTML}{F4D94B}
\definecolor{MarginError}{HTML}{F28A18}
\colorlet{Endpoint}{MarginError!45!white}
\colorlet{Width}{MarginError!45!black}
\definecolor{Coverage}{HTML}{BC9DE0}
\tikzset{ge marker/.style={rectangle,rounded corners=1mm,draw=none,
  fill=#1,fill opacity=.32,text opacity=1,text=GEInk,
  inner xsep=.6mm,inner ysep=.3mm,outer sep=0pt}}
\makeatletter
\newcommand{\GEHighlightMath}[2]{%
  \tikz[baseline=(gehl.base)]\node[ge marker=\GECurrentRole]
    (gehl) {$\m@th#1#2$};}
\newcommand{\GEHighlight}[2]{%
  \ifGTenHighlights
  \begingroup\def\GECurrentRole{#1}%
  \ifmmode\mathpalette\GEHighlightMath{#2}%
  \else\tikz[baseline=(gehl.base)]\node[ge marker=#1](gehl){#2};\fi
  \endgroup
  \else#2\fi}
\makeatother
\newif\ifGTenHighlights
\GTenHighlightstrue
% \GTenHighlightsfalse % Uncomment to disable instructional highlights.
% Keep rows, rules, and table structure outside this space-reserving command.
\newcommand{\GERevealCell}[2]{\leavevmode\visible<#1->{#2}}


\newcommand{\CtwoMean}[1]{\GEHighlight{PopulationMean}{#1}}
\newcommand{\CtwoPoint}[1]{\GEHighlight{SampleMean}{#1}}
\newcommand{\CtwoSD}[1]{\GEHighlight{PopulationSD}{#1}}
\newcommand{\CtwoSize}[1]{\GEHighlight{SampleSize}{#1}}
\newcommand{\CtwoConfidence}[1]{\GEHighlight{Confidence}{#1}}
\newcommand{\CtwoSignificance}[1]{\GEHighlight{Confidence}{#1}}
\newcommand{\CtwoCritical}[1]{\GEHighlight{Critical}{#1}}
\newcommand{\CtwoSE}[1]{\GEHighlight{StandardError}{#1}}
\newcommand{\CtwoMoE}[1]{\GEHighlight{MarginError}{#1}}
\newcommand{\CtwoEndpoint}[1]{\GEHighlight{Endpoint}{#1}}
\newcommand{\CtwoWidth}[1]{\GEHighlight{Width}{#1}}
\newcommand{\CtwoCoverage}[1]{\GEHighlight{Coverage}{#1}}
\newcommand{\GELegendCell}[2]{\GERevealCell{#1}{#2}}
\newcommand{\GEFullLegendCell}[2]{\GERevealCell{1}{#2}}
\newcommand{\GEGeneralNotationLegend}{%
\begingroup\small\setlength{\parskip}{0pt}
\begin{tabularx}{\linewidth}{@{}>{\raggedright\arraybackslash}p{.46\linewidth}>{\raggedright\arraybackslash}X@{}}
\toprule
\textbf{Terms and symbols} & \textbf{Contextual meaning and units}\\\midrule
\GELegendCell{1}{\CtwoMean{Population mean}, $\CtwoMean{\mu}$} & \GELegendCell{1}{Fixed parameter; minutes or thousand COP.}\\
\GELegendCell{1}{\CtwoPoint{Point estimate}, $\CtwoPoint{\bar{x}}$} & \GELegendCell{1}{Sample mean / interval centre; units of $\mu$.}\\
\GELegendCell{1}{\CtwoSD{Population SD $\sigma$}; \CtwoSize{sample size $n$}} & \GELegendCell{1}{Individual spread (units of $\mu$); sample count.}\\
\GELegendCell{2}{\CtwoSE{Standard error}, $\CtwoSE{\mathrm{SE}}$} & \GELegendCell{2}{Sample-mean spread (or estimate); units of $\mu$.}\\
\GELegendCell{2}{\CtwoCritical{Critical value}, $\CtwoCritical{z^*}$} & \GELegendCell{2}{Standard normal multiplier set by $C$; unitless.}\\
\GELegendCell{2}{\CtwoMoE{Margin of error}, $\CtwoMoE{\mathrm{MoE}}$} & \GELegendCell{2}{Centre-to-bound distance $z^*\mathrm{SE}$; units of $\mu$.}\\
\GELegendCell{2}{\CtwoEndpoint{Endpoints}; \CtwoWidth{full width}} & \GELegendCell{2}{Bounds $\bar{x}\pm\mathrm{MoE}$; span $2\mathrm{MoE}$; units of $\mu$.}\\
\GELegendCell{3}{\CtwoConfidence{Confidence $C$}; \CtwoSignificance{significance $\alpha$}} & \GELegendCell{3}{Long-run $C$; $\alpha=1-C$; tails $\alpha/2$ (no units).}\\
\GELegendCell{3}{\CtwoCoverage{Observed coverage}, rate} & \GELegendCell{3}{Finite-run coverage (\%); can differ from $C$.}\\
\bottomrule
\end{tabularx}\par\endgroup}
\title{Point estimates and confidence levels}
\subtitle{Grade 11 · Term 2 · Lesson 1\\Global Economics · 2026--2027}
\author{}\date{}
\newcommand{\GEOfficialCriterion}{\textbf{C2} Compares point and interval estimation and interprets confidence levels through their long-run coverage meaning, distinguishing this from the probability that a particular calculated interval contains the parameter.}
\setbeamertemplate{title page}{\raggedright\vfill
{\color{GETeal}\usebeamerfont{subtitle}\insertsubtitle\par}\medskip
{\color{GEInk}\usebeamerfont{title}\inserttitle\par}\bigskip
{\normalsize\GEOfficialCriterion\par}\vfill}
\begin{document}
\section{Introduction}
\begin{frame}\titlepage\end{frame}
\begin{frame}[t]{General notation table legend}\GEGeneralNotationLegend\end{frame}

% Introduction, subsec:c2-cycling-example: motivation precedes explanation.
\begin{frame}[t]{Is the average cycling trip under 42 minutes?}
A planner surveys an independent random sample of cyclists making weekly trips to the market.

Assume a normal sampling model and known population standard deviation.
\par\medskip
\textbf{Predict:} Would a sample mean of \CtwoPoint{40 minutes} alone settle whether the population mean is under 42 minutes?
\par\medskip
\visible<2->{The sample mean is a \CtwoPoint{point estimate}. An interval gives a range of plausible values for the population mean under the model.}
\end{frame}

\begin{frame}[t]{Key terms and highlight guide}
A \CtwoMean{population parameter} describes the population; here, the \CtwoMean{population mean} is fixed.

A \CtwoPoint{point estimate} uses a sample statistic to estimate that parameter. An \textbf{interval estimate} gives a range of plausible parameter values under the model.
\medskip
\visible<2->{The \CtwoConfidence{confidence level} describes the procedure's long-run coverage. Its complement is the \CtwoSignificance{significance level}:}
\[
\visible<2->{\CtwoSignificance{\alpha}=1-\CtwoConfidence{C}\qquad(C\text{ expressed as a decimal}).}
\]
\visible<3->{Greens: sample quantities. Reds: population quantities. Blues: confidence/significance and critical value. Yellow: SE. Orange: MoE, lighter endpoints, darker width. Names and symbols distinguish roles without color.}
\end{frame}

% All generic numbers on these two frames belong to the introduction's guide.
\begin{frame}[t]{Reading the guide: spread and interval}
The guide's illustrative values use generic \textbf{units}.
\[
\CtwoPoint{\bar{x}}=\CtwoPoint{60\,\text{units}},\quad
\CtwoSD{\sigma}=\CtwoSD{20\,\text{units}},\quad \CtwoSize{n}=\CtwoSize{100}.
\]
\textbf{Interpret:} Which spread describes individuals, and which describes sample means?
\[
\visible<2->{\CtwoSE{\mathrm{SE}}=\frac{\CtwoSD{20}}{\sqrt{\CtwoSize{100}}}=\CtwoSE{2\,\text{units}}.}
\]
\visible<2->{$\sigma$ describes individual observations; SE describes sample means across repeated samples, or estimates that spread.}
\[
\visible<3->{\CtwoMoE{\mathrm{MoE}}=\CtwoCritical{1.96}(\CtwoSE{2})=\CtwoMoE{3.92\,\text{units}}.}
\]
\end{frame}

\begin{frame}[t]{Reading the guide: bounds and coverage}
\textbf{Interpret:} What do $\CtwoPoint{60}\pm\CtwoMoE{3.92}$ and $2(\CtwoMoE{3.92})$ describe?
\[
\visible<2->{[\CtwoEndpoint{56.08},\CtwoEndpoint{63.92}]\,\text{units},\qquad
\CtwoWidth{\text{width}}=\CtwoWidth{7.84\,\text{units}}.}
\]
\visible<2->{The endpoints bound the interval; full width is twice its margin of error.}
\medskip
\visible<3->{For $\CtwoConfidence{C}=\CtwoConfidence{0.95}$, $\CtwoSignificance{\alpha}=\CtwoSignificance{0.05}$ and each tail has area $\CtwoSignificance{0.025}$. The critical bounds $\pm\CtwoCritical{1.96}$ enclose central normal area \CtwoConfidence{0.95}.}
\medskip
\visible<4->{If 4 of 5 intervals cover, \CtwoCoverage{observed coverage} is $4/5=\CtwoCoverage{80\%}$. Each interval covers or misses. This finite-run rate can differ from $C$; it does not set $\alpha$.}
\end{frame}

\begin{frame}[t]{Cycling study: organize the information}
\textbf{Task 1:} Identify population, point estimate, sample design, and procedure settings; calculate SE.
\medskip
\begin{tabularx}{\linewidth}{@{}>{\raggedright\arraybackslash}p{.30\linewidth}X@{}}\toprule
\textbf{Role} & \textbf{Original study}\\\midrule
Population parameter & \GERevealCell{2}{All weekly cycling trips to the market; unknown $\CtwoMean{\mu}$ in \CtwoMean{minutes}.}\\
Sample statistic & \GERevealCell{2}{$\CtwoPoint{\bar{x}}=\CtwoPoint{40\,\text{minutes}}$, a point estimate of $\mu$.}\\
Sample design & \GERevealCell{3}{$\CtwoSize{n}=\CtwoSize{49\,\text{trips}}$.}\\
Known population SD & \GERevealCell{3}{$\CtwoSD{\sigma}\approx\CtwoSD{14.29\,\text{minutes}}$.}\\
Procedure settings & \GERevealCell{4}{$\CtwoConfidence{C}=\CtwoConfidence{0.95}$; $\CtwoSignificance{\alpha}=1-\CtwoConfidence{0.95}=\CtwoSignificance{0.05}$;\par $\CtwoSignificance{\alpha}/2=\CtwoSignificance{0.025}$; $\CtwoCritical{z^*}=\CtwoCritical{1.960}$.}\\\bottomrule
\end{tabularx}
\end{frame}

\begin{frame}[t]{Cycling study: calculate and interpret}
\textbf{Tasks 1--2:} Calculate SE, MoE, and the interval. Does it support a mean under 42 minutes?

Population SDs shown with $\approx$ are rounded to two decimals; calculations retain their unrounded values.
\[
\begin{aligned}
\visible<2->{\CtwoSE{\mathrm{SE}}&=\CtwoSD{\sigma}/\sqrt{\CtwoSize{49}}\approx\frac{\CtwoSD{14.29}}{\sqrt{\CtwoSize{49}}}\approx\CtwoSE{100/49\,\text{minutes}}.}\\[5pt]
\visible<3->{\CtwoMoE{\mathrm{MoE}}&=\CtwoCritical{1.960}(\CtwoSE{100/49})=\CtwoMoE{4\,\text{minutes}}.}\\[5pt]
\visible<4->{[\CtwoPoint{40}-\CtwoMoE{4},\CtwoPoint{40}+\CtwoMoE{4}]&=[\CtwoEndpoint{36},\CtwoEndpoint{44}]\,\CtwoEndpoint{\text{minutes}}.}
\end{aligned}
\]
\visible<5->{These are plausible values for the population mean under the model. The interval includes 42 and larger values, so it does not support concluding that the mean is under 42 minutes.}
\end{frame}

\begin{frame}[t]{Cycling simulation: what fraction covers?}
\textbf{Task 3:} Separately, a hypothetical simulation uses known fixed mean $\CtwoMean{\mu}=\CtwoMean{41\,\text{minutes}}$ and the same $\sigma$, $n$, and 95\% procedure.

Of 50 fresh samples, 3 give intervals that miss 41. Calculate observed coverage.
\[
\visible<2->{50-3=47\text{ covering intervals},\qquad
\frac{50-3}{50}\times100\%=\CtwoCoverage{94\%}.}
\]
\visible<3->{These are new simulated samples, not the original survey. The finite-run \CtwoCoverage{94\%} differs from nominal \CtwoConfidence{95\%}; it does not change $\CtwoSignificance{\alpha}=\CtwoSignificance{0.05}$.}
\medskip
\visible<4->{The nominal level describes about \CtwoConfidence{95\%} coverage over many repetitions under the model.}
\end{frame}

\begin{frame}[t]{What does confidence describe?}
\textbf{Interpret:} In the original cycling study, is confidence a percentage of individual trip times within the interval?
\par\medskip
\visible<2->{For the original study, the fixed mean remains unknown. A realized interval either contains it or does not.}
\par\medskip
\visible<3->{Confidence does not give the percentage of individual trip times within the interval. It describes the interval procedure's long-run coverage.}
\end{frame}

\section{Worked example}
% prob:c2-2, complete setup, all three tasks and matching solutions.
\begin{frame}[t]{Two reports on household energy bills}
A utility takes a random sample of $\CtwoSize{n}=\CtwoSize{100}$ household accounts from one billing month. Both reports use this sample:
\[
\CtwoPoint{\bar{x}}=\CtwoPoint{120\,\text{thousand COP}},\qquad
\CtwoSD{\sigma}=\CtwoSD{38\,\text{thousand COP}}.
\]
Assume independent observations and a normal sampling model; population SD is known.
\medskip
\begin{tabularx}{\linewidth}{@{}XXX@{}}\toprule
Nominal confidence & Critical value $z^*$ & Width (thousand COP)\\\midrule
\CtwoConfidence{90\%} & \CtwoCritical{1.645} & To calculate\\
\CtwoConfidence{99\%} & \CtwoCritical{2.576} & To calculate\\\bottomrule
\end{tabularx}
\end{frame}

\begin{frame}[t]{Task 1: organize the shared information}
Identify population parameters, sample statistic, and design information. What is shared by both reports?
\medskip
\begin{tabularx}{\linewidth}{@{}>{\raggedright\arraybackslash}p{.28\linewidth}X@{}}\toprule
\textbf{Role} & \textbf{Information and meaning}\\\midrule
Population parameter & \GERevealCell{2}{$\CtwoMean{\mu}$: fixed, unknown mean energy bill for all the utility's household accounts in that month, in \CtwoMean{thousand COP}.}\\
Point estimate & \GERevealCell{3}{$\CtwoPoint{\bar{x}}=\CtwoPoint{120\,\text{thousand COP}}$: a sample statistic giving a single estimate of the monthly mean.}\\
Sample design & \GERevealCell{4}{$\CtwoSize{n}=\CtwoSize{100\,\text{accounts}}$; both reports use the same sample.}\\
Population SD & \GERevealCell{4}{$\CtwoSD{\sigma}=\CtwoSD{38\,\text{thousand COP}}$; shared known parameter.}\\\bottomrule
\end{tabularx}
\end{frame}

\begin{frame}[t]{Task 1: procedure settings and SE}
Identify $\alpha=1-C$ and each tail area $\alpha/2$. Calculate the shared SE.
\medskip
\begin{tabularx}{\linewidth}{@{}XXXX@{}}\toprule
$\CtwoConfidence{C}$ & $\CtwoSignificance{\alpha}$ & $\CtwoSignificance{\alpha}/2$ & $\CtwoCritical{z^*}$\\\midrule
\CtwoConfidence{0.90} & \GERevealCell{2}{\CtwoSignificance{0.10}} & \GERevealCell{2}{\CtwoSignificance{0.05}} & \CtwoCritical{1.645}\\
\CtwoConfidence{0.99} & \GERevealCell{3}{\CtwoSignificance{0.01}} & \GERevealCell{3}{\CtwoSignificance{0.005}} & \CtwoCritical{2.576}\\\bottomrule
\end{tabularx}
\[
\visible<4->{\CtwoSE{\mathrm{SE}}=\frac{\CtwoSD{\sigma}}{\sqrt{\CtwoSize{n}}}
=\frac{\CtwoSD{38}}{\sqrt{\CtwoSize{100}}}=\CtwoSE{3.8\,\text{thousand COP}}.}
\]
\visible<4->{Both intervals use the same sample and standard error. Their nominal levels, significance levels, tail areas, and critical values differ.}
\end{frame}

\begin{frame}[t]{Task 2: calculate and report both intervals}
For each procedure, calculate MoE and the interval; briefly interpret it in context.
\[
\CtwoMoE{\mathrm{MoE}}=\CtwoCritical{z^*}\times\CtwoSE{\mathrm{SE}},\qquad
[\CtwoPoint{\bar{x}}-\CtwoMoE{\mathrm{MoE}},\CtwoPoint{\bar{x}}+\CtwoMoE{\mathrm{MoE}}].
\]
State unrounded endpoints and widths. Then round endpoints to the nearest thousand COP and complete the width table.
\medskip
The confidence levels refer to the procedures \textbf{before rounding}.
\end{frame}

\begin{frame}[t]{Task 2: the 90\% interval}
Calculate MoE and endpoints using $\CtwoPoint{\bar{x}}=\CtwoPoint{120}$ and $\CtwoSE{\mathrm{SE}}=\CtwoSE{3.8}$, in thousand COP.
\[
\begin{aligned}
\visible<2->{\CtwoMoE{\mathrm{MoE}}&=\CtwoCritical{1.645}(\CtwoSE{3.8})=\CtwoMoE{6.251\,\text{thousand COP}}.}\\[6pt]
\visible<3->{\CtwoPoint{120}\pm\CtwoMoE{6.251}&\longrightarrow[\CtwoEndpoint{113.749},\CtwoEndpoint{126.251}].}\\[6pt]
\visible<4->{\text{Rounded interval}&=[\CtwoEndpoint{114},\CtwoEndpoint{126}]\,\CtwoEndpoint{\text{thousand COP}}.}
\end{aligned}
\]
\visible<5->{The narrower reported range of plausible values for the mean household bill under the model.}
\end{frame}

\begin{frame}[t]{Task 2: the 99\% interval}
Calculate MoE and endpoints using the same sample and SE, in thousand COP.
\[
\begin{aligned}
\visible<2->{\CtwoMoE{\mathrm{MoE}}&=\CtwoCritical{2.576}(\CtwoSE{3.8})=\CtwoMoE{9.7888\,\text{thousand COP}}.}\\[6pt]
\visible<3->{\CtwoPoint{120}\pm\CtwoMoE{9.7888}&\longrightarrow[\CtwoEndpoint{110.2112},\CtwoEndpoint{129.7888}].}\\[6pt]
\visible<4->{\text{Rounded interval}&=[\CtwoEndpoint{110},\CtwoEndpoint{130}]\,\CtwoEndpoint{\text{thousand COP}}.}
\end{aligned}
\]
\visible<5->{The wider reported range of plausible values for the same mean under the model.}
\end{frame}

\begin{frame}[t]{Task 2: compare widths}
How do unrounded and reported widths compare? All widths below are in thousand COP.
\medskip
\begin{tabularx}{\linewidth}{@{}>{\raggedright\arraybackslash}p{.13\linewidth}XX@{}}\toprule
Level & Unrounded width & Reported width\\\midrule
\CtwoConfidence{90\%} & \GERevealCell{2}{$2(\CtwoMoE{6.251})=\CtwoWidth{12.502}$} & \GERevealCell{3}{$\CtwoEndpoint{126}-\CtwoEndpoint{114}=\CtwoWidth{12}$}\\
\CtwoConfidence{99\%} & \GERevealCell{2}{$2(\CtwoMoE{9.7888})=\CtwoWidth{19.5776}$} & \GERevealCell{3}{$\CtwoEndpoint{130}-\CtwoEndpoint{110}=\CtwoWidth{20}$}\\\bottomrule
\end{tabularx}
\medskip
\visible<4->{With the same SE, the larger critical value for \CtwoConfidence{99\%} gives a larger MoE and puts the endpoints farther from \CtwoPoint{120}.}
\end{frame}

\begin{frame}[t]{Task 3: separate simulation evidence}
In a separate, explicitly hypothetical simulation, a normal model has known fixed mean $\CtwoMean{\mu}=\CtwoMean{120\,\text{thousand COP}}$ and the same $\CtwoSD{\sigma}=\CtwoSD{38}$.

From each of 100 fresh samples of size $\CtwoSize{n}=\CtwoSize{100}$, the utility constructs both an unrounded 90\% interval and an unrounded 99\% interval using the stated critical values.
\medskip
Of the 100 unrounded 90\% intervals, \textbf{11 miss} the mean. Of the 100 unrounded 99\% intervals, \textbf{2 miss} it.

Coverage is evaluated \textbf{before rounding}, separately for each level. These simulated samples are not the original household sample.
\end{frame}

\begin{frame}[t]{Task 3: choose an interpretation (A--B)}
Which option gives observed coverage for the 90\% and 99\% procedures, in that order, correctly compares the widths, and interprets 99\% confidence?
\par\medskip
\textbf{A.} The observed rates are 89\% and 98\%, respectively; the 99\% interval is wider and contains 99\% of individual bills.
\par\medskip
\textbf{B.} The observed rates are 90\% and 99\%, respectively; each finite run must match its nominal confidence level.
\end{frame}

\begin{frame}[t]{Task 3: choose an interpretation (C--D)}
Which option gives the rates, compares widths, and interprets 99\% confidence?
\par\medskip
\textbf{C.} The observed rates are 89\% and 98\%, respectively; the wider reported interval has a 99\% probability of containing the fixed mean.
\par\medskip
\textbf{D.} The observed rates are 89\% and 98\%, respectively. The 99\% interval is wider; its unrounded procedure covers the fixed mean about 99\% of the time in the long run under the model, without specifying a percentage of individual bills inside it.
\end{frame}

\begin{frame}[t]{Task 3: calculate observed coverage}
How many intervals cover for each procedure? What is each denominator?
\[
\begin{aligned}
\visible<2->{90\%:\quad 100-11&=89\text{ covering intervals},}\\[5pt]
\visible<3->{\frac{89}{100}\times100\%&=\CtwoCoverage{89\%}.}\\[8pt]
\visible<4->{99\%:\quad 100-2&=98\text{ covering intervals},}\\[5pt]
\visible<5->{\frac{98}{100}\times100\%&=\CtwoCoverage{98\%}.}
\end{aligned}
\]
\visible<6->{Each denominator is 100, not 200. Both counts concern unrounded intervals. The finite-run rates need not equal the nominal levels.}
\end{frame}

\begin{frame}[t]{Task 3: explain the choice}
Which option remains consistent with the calculations and the meaning of confidence?
\par\medskip
\visible<2->{\textbf{Correct option D.} Under the model assumptions, about \CtwoConfidence{90\%} and \CtwoConfidence{99\%} of the respective unrounded intervals would cover the fixed mean over repeated samples.}
\par\medskip
\visible<3->{Each calculated interval either \CtwoCoverage{covers the mean} or \CtwoCoverage{misses it}. Confidence does not assign a 99\% probability to that fixed interval containing the fixed mean.}
\end{frame}

\begin{frame}[t]{Task 3: the mean and individual bills}
Why does the claim about 99\% of individual bills fail?
\par\medskip
\visible<2->{Coverage concerns the \CtwoMean{population mean}.}
\par\medskip
\visible<3->{The proportion of individual bills from \CtwoEndpoint{110} to \CtwoEndpoint{130 thousand COP} depends on the distribution of individual bills.}
\end{frame}

\section{Summary}
\begin{frame}[t]{Point, interval, and confidence}
\begin{columns}[T,onlytextwidth]
\begin{column}{.485\textwidth}
\textbf{Estimation}\par
A \CtwoPoint{point estimate} gives one value; an interval gives plausible values for the fixed \CtwoMean{population mean} under the model.
\par\medskip
With the same SE, the larger critical value gives a larger MoE and a wider interval.
\end{column}
\begin{column}{.485\textwidth}
\textbf{Coverage}\par
\CtwoConfidence{Nominal confidence} describes long-run coverage. \CtwoCoverage{Observed coverage} is a finite-run rate.
\par\medskip
A realized interval covers or misses the fixed mean; confidence does not describe the percentage of individual observations within it.
\end{column}
\end{columns}
\end{frame}
\section{Reminder}
\begin{frame}[t]{General notation table legend}
\let\GELegendCell\GEFullLegendCell
\GEGeneralNotationLegend
\end{frame}
\end{document}
